\head{第二十一版整合核查：第13节的实际范围}
\audit{V21-INT01}{接续基线与新回校范围}{以可读源码和底本为依据，不从旧聊天推定完成度}
本轮基线为公开提交\texttt{99cd048}中的第二十版工程，其新回校点确实止于第一篇第12节。既有第13节已经有中文译文；本版将该段与法稿物理74--82页、英编稿49--53页逐段对照，形成三份新正文模块。这里记作已有译文的法英回校、恢复与订正，不计成从未译过的新增整节。

\begin{longtable}{@{}p{.18\linewidth}p{.21\linewidth}p{.49\linewidth}@{}}
\toprule 模块 & 底本位置 & 精确内容边界\\\midrule\endhead
初段 & 法74--77；英49--51首段 & 节标题、13.1、13.2、13.3陈述，以及初等曲线族构造，止于所求邻域。\\
截面证明 & 法78--79及80页顶部；英51中--52页附记末 & 引理13.3证明、固定点截面、归纳、对偶论证及正维中心附记。\\
五项注 & 法80中--82末；英52末--53 & 13.4(a)--(e)，包括有限域、弱化定义、有限点、一般基底及解析拓扑。\\
\bottomrule
\end{longtable}

法83页及英53页末至54页只用于确认下一回校起点：第14节标题在英文53页末，正文在54页。没有将第14节算作本轮已经完成的回校内容。法文纸面13.4诸项的次序与英文编排不同，已在本版注释中明说，五项内容均保留。

\audit{V21-INT02}{数学订正的性质}{原错误、英文误读、补证和未决事项分别列明}
\lead{可以直接据法稿恢复的内容。}
包含链$X\supset Y\supset Z$、$C\subset L_0$、一般点下标、先超平面后子空间的归纳顺序、表示“问题”的Pb、条件连接词、$\Phi(\bar k)$及$U^{\mathrm{an}}$等，法稿有明确依据。相关旧报告曾依据残缺英文归责于原作者，本版撤销或限定这些归责。

\lead{法稿确有的错误或遗漏。}
定义13.1中的紧化态射靶、邻域所在的参数空间、被删去的边界字母、声称包含整个截面而未删边界、一般纤维的基域，以及13.4(c)的一次余维数均已分别订正。横截中心及共同处理非有理点所需的步骤，不以“显然”掩盖，报告另列完整补证。

\lead{反例的范围。}
不完美域上的正则非光滑射影曲线说明13.2的完美性不能无条件删除。三次Veronese曲线的切线平面束否定正维固定中心附记中不附限制的乐观猜测；原作者只是提问，不把它说成已证明的错误定理。非平凡有限可分扩张的谱说明13.4(b)仅删连通条件不足以包含零维。三个反例不互相替代，也不否定原13.2在无限完美域、几何不可约条件下的结论。

\lead{补证的范围。}
一阶截面计算给出次数至少2的Veronese充分版本；有限点归纳明确携带全部几何点，并用有限平展集合下降共同参数条件。连通性所用的射影Bertini定理有独立文献支点，不冒称本版重新证明了一般连通性定理。解析拓扑部分以保持边界的平行移动建立欧氏普遍覆盖，而不从可缩性倒推出同胚结论；正文只保留原稿的同胚表述，没有改成双全纯。

\lead{仍然未决。}
原8.7(c)、8.8的缺续没有恢复，特征零不改给定嵌入的一般固定点论证没有在本版另行全面重证。13.4(a)没有消去有限扩域，13.4(d)在任意基预概形上的全部推广仍没有完成核查。不可辨的密集覆打不据数学需要补造。其他卷、其他章节及历史报告不因本轮局部审校而全部获得新的数学认证。

\audit{V21-INT03}{源码保留、构建与交接}{技术记录与数学结论分开}
第二十版的正文、报告和书末TeX文件作为不可变继承对象保留。原第13节文件及其33,351字节副本均保存；当前输入链以三份新模块替换这一个旧入口，避免重复排印。既有Dot69--108共40个模块仍各输入一次，Dot后续原包和既有档案未改写。旧报告全部接在本版核查之后，并标明各自历史版次。

使用真正的\texttt{ctexbook}默认字体、字号和版心，继承导言文件不作全局数学宏重定义。新增正文仍用TeX排公式，单文件通过展开当前输入链生成。工程中的结构核验、四份最终XeLaTeX日志、模块与单文件逐页文字/尺寸/像素比较，以及绑定成品散列的实际读图记录，分别用于检查不同问题；具体页数和检查范围见本工程的实际记录，不能把样页目视写成全书每页人工审核。

当前接续文档明确记录：下一是第一篇第14节，法83页，英53页末标题及54页正文。正式第IV章主体仍是另一个未完成分支，从其第一分册物理221页提要及后续章首开始。所有书籍尚未完结；既有历史容器中7个早期原ZIP不可恢复的事实继续保留，不能因为普通成员文件仍在就声称原容器也已齐全。
